\documentclass[11pt]{article}

\usepackage[margin=1in]{geometry}
\usepackage{amsmath}
\usepackage{amssymb}
\usepackage{booktabs}
\usepackage{hyperref}

\title{Quantitative Pathway for Protocell Formation Phases}
\author{LIFEORIG working note}
\date{\today}

\begin{document}
\maketitle

\section{Purpose}

This note defines a quantitative pathway for modeling different physical
stages of protocell formation inside the LIFEORIG framework. The aim is not to
declare one origin-of-life pathway correct. The aim is to define computable
states, transition criteria, and equations that can later be coupled to the
atmosphere, hydrosphere, ocean chemistry, geophysical interior, and prebiotic
reaction network.

The phases considered here are:
\begin{enumerate}
    \item porous mineral or ice compartments;
    \item solid surface compartments and adsorbed chemical networks;
    \item liquid-liquid coacervate compartments;
    \item amphiphilic vesicles and true protocell-like compartments.
\end{enumerate}

The working philosophy is:
\begin{equation}
    \mathrm{planetary\ environment}
    \rightarrow
    \mathrm{hydrosphere/geochemistry}
    \rightarrow
    \mathrm{chemical\ network}
    \rightarrow
    \mathrm{compartment\ physics}
    \rightarrow
    \mathrm{growth/division/selection}.
\end{equation}

\section{Environmental Inputs}

The protocell module should not invent its own environment. It should receive
boundary conditions from the atmosphere, hydrosphere, and geochemistry solvers.

\subsection{Atmosphere and hydrosphere inputs}

Required environmental variables:
\begin{align}
    T_s       &\quad \mathrm{surface\ temperature}\ [\mathrm{K}], \\
    P_s       &\quad \mathrm{surface\ pressure}\ [\mathrm{Pa}], \\
    F_{\mathrm{UV}}(\lambda,z)
              &\quad \mathrm{local\ UV\ photon\ flux}, \\
    \tau_{\mathrm{UV}}(\lambda,z)
              &\quad \mathrm{UV\ optical\ depth}, \\
    M_w        &\quad \mathrm{available\ liquid/ice\ water\ mass}, \\
    d_w        &\quad \mathrm{surface\ water\ depth\ or\ pore\ water\ depth}, \\
    a_w        &\quad \mathrm{water\ activity}, \\
    \Phi_{\mathrm{wetdry}}
              &\quad \mathrm{wet/dry\ cycling\ strength}, \\
    \Phi_{\mathrm{freeze}}
              &\quad \mathrm{freeze/thaw\ concentration\ factor}.
\end{align}

For a subsurface ocean, use the ocean pressure and temperature at the relevant
depth:
\begin{align}
    T_o &= T(r_{\mathrm{ocean}}), \\
    P_o &= P(r_{\mathrm{ocean}}).
\end{align}

For a surface pond or volcanic-rock environment, use local surface values:
\begin{align}
    T_{\mathrm{env}} &= T_s, \\
    P_{\mathrm{env}} &= P_s.
\end{align}

\subsection{Geochemical inputs}

The geochemistry module should provide:
\begin{align}
    \mathrm{pH},\quad
    I,\quad
    \mathrm{Eh},\quad
    [\mathrm{Fe^{2+}}],\quad
    [\mathrm{Fe^{3+}}],\quad
    [\mathrm{HS^-}],\quad
    [\mathrm{SO_4^{2-}}],\quad
    [\mathrm{PO_4}],\quad
    [\mathrm{Mg^{2+}}],
\end{align}
plus mineral availability:
\begin{equation}
    S_m(X) \quad \mathrm{reactive\ surface\ area\ of\ mineral\ } X.
\end{equation}

The relevant mineral classes are:
\begin{equation}
    X \in
    \{\mathrm{clay},\mathrm{FeS},\mathrm{NiS},\mathrm{FeOOH},
    \mathrm{silica},\mathrm{phosphate\ surface},\mathrm{ice}\}.
\end{equation}

\subsection{Chemical network inputs}

The reaction network should provide concentrations or activities of:
\begin{align}
    &\mathrm{HCN},\ \mathrm{H_2CO},\ \mathrm{NH_3},\ \mathrm{CO},\
      \mathrm{CO_2},\ \mathrm{H_2},\ \mathrm{H_2S},\ \mathrm{CH_3SH}, \\
    &\mathrm{formamide},\ \mathrm{formate},\ \mathrm{acetate},\
      \mathrm{glyoxylate},\ \mathrm{glycolaldehyde},\
      \mathrm{pyruvate}, \\
    &\mathrm{glycine},\ \mathrm{alanine},\ \mathrm{aspartate},\
      \mathrm{glutamate}, \\
    &\mathrm{thioacetate},\ \mathrm{acetyl\ phosphate},\
      \mathrm{aminoacyl\ thioesters}, \\
    &\mathrm{peptide\ oligomers},\ \mathrm{hetero\ oligomers},\
      \mathrm{amphiphiles}.
\end{align}

\section{Central Mechanism: Wet-Dry Concentration and Chemical Trapping}

The difficult step is not only producing organics. The difficult step is
keeping their concentration high enough, for long enough, to build boundaries.
Therefore protocell formation should be modeled as a transient environmental
cycle, not as a simple average-abundance threshold.

The core sequence is:
\begin{equation}
    \mathrm{wet/dry\ cycle}
    \rightarrow
    \mathrm{concentration\ pulse}
    \rightarrow
    \mathrm{amphiphile\ synthesis/assembly}
    \rightarrow
    \mathrm{boundary\ closure}
    \rightarrow
    \mathrm{chemical\ trapping}
    \rightarrow
    \mathrm{outside\ dilution}.
\end{equation}

\subsection{Water-depth controlled concentration}

Let $h(t)$ be the local water depth in a pond, pore, film, or wet mineral
surface. If the number of molecules $n_i$ is approximately conserved during
drying and the wetted area is $A_{\mathrm{pond}}$, then
\begin{equation}
    c_i(t)
    =
    \frac{n_i(t)}
    {A_{\mathrm{pond}}h(t)}.
\end{equation}

If chemistry is slow compared with the drying concentration step, then
\begin{equation}
    c_i(t) \propto \frac{1}{h(t)}.
\end{equation}

This is important because many prebiotic reactions are second order or higher.
For example, if an amphiphile precursor $L$ forms by
\begin{equation}
    A+B\rightarrow L,
\end{equation}
then
\begin{equation}
    \frac{dc_L}{dt}
    =
    k_L(T,\mathrm{pH})c_Ac_B
    -
    k_{\mathrm{loss}}c_L.
\end{equation}

If drying increases both $c_A$ and $c_B$ by a factor $\Gamma$, the production
term increases as $\Gamma^2$.

\subsection{Explicit wet-dry concentration term}

A minimal local environmental equation is:
\begin{equation}
    \frac{dc_i^{\mathrm{ext}}}{dt}
    =
    R_i^{\mathrm{env}}(\mathbf{c},T,\mathrm{pH})
    +
    S_i^{\mathrm{geo}}
    -
    k_{\mathrm{wash}}W(t)c_i^{\mathrm{ext}}
    -
    \frac{c_i^{\mathrm{ext}}}{h}
    \frac{dh}{dt}.
\end{equation}

The last term is the concentration/dilution term. During drying,
$dh/dt<0$, so the term is positive. During rehydration, $dh/dt>0$, so the
term is negative.

A simple periodic water depth can be:
\begin{equation}
    h(t)
    =
    h_0
    +
    \Delta h
    \cos\left(\frac{2\pi t}{\tau_{\mathrm{cyc}}}\right),
\end{equation}
with clipping if needed to avoid negative water depth.

The concentration factor relative to a reference depth $h_0$ is:
\begin{equation}
    \Gamma(t) =
    \frac{h_0}{h(t)}.
\end{equation}

\subsection{Assembly exposure}

The relevant condition is not merely instantaneous threshold crossing. The
system must remain above threshold long enough for assembly.

Let $c_L$ be the total concentration of membrane-forming amphiphiles:
\begin{equation}
    c_L(r,t) =
    \sum_{\ell} w_\ell c_{L_\ell}(r,t).
\end{equation}

Define the accumulated amphiphile assembly exposure:
\begin{equation}
    Q_L(r,t)
    =
    \int_0^t
    dt'
    \max\left[
    0,\,
    \frac{c_L(r,t')-c_L^{\mathrm{crit}}}
    {c_L^{\mathrm{crit}}}
    \right].
\end{equation}

Vesicle formation is allowed when:
\begin{equation}
    Q_L(r,t) > Q_L^{\mathrm{crit}}.
\end{equation}

Equivalently, if the concentration is above threshold over a dry interval
$\Delta t_{\mathrm{asm}}$:
\begin{equation}
    \Delta t_{\mathrm{asm}}
    =
    \int dt\,
    \Theta\left[c_L(t)-c_L^{\mathrm{crit}}\right],
\end{equation}
then assembly requires:
\begin{equation}
    \Delta t_{\mathrm{asm}} \gtrsim \tau_{\mathrm{asm}}.
\end{equation}

\subsection{Boundary closure traps the chemical state}

If a vesicle forms at $t=t_{\mathrm{form}}$, its initial contents should be
sampled from the environmental concentration at that time:
\begin{equation}
    X_i^k(t_{\mathrm{form}})
    \sim
    \mathrm{Poisson}
    \left[
    c_i^{\mathrm{ext}}(r_k,t_{\mathrm{form}})V_k^0
    \right].
\end{equation}

This is the chemical memory of the dry phase. After rehydration, $h(t)$ rises
and $c_i^{\mathrm{ext}}$ decreases, but the internal contents do not instantly
dilute because a boundary exists.

\subsection{Leakage after rehydration}

For a vesicle, passive exchange is:
\begin{equation}
    J_i^k =
    A_k P_i
    \left[
    c_i^{\mathrm{ext}}(r_k,t)
    -
    c_i^{k,\mathrm{int}}(t)
    \right].
\end{equation}

After rehydration,
\begin{equation}
    c_i^{\mathrm{ext}} < c_i^{k,\mathrm{int}},
\end{equation}
so $J_i^k<0$ and molecules leak out.

The leakage timescale is approximately:
\begin{equation}
    \tau_{i,\mathrm{leak}}
    =
    \frac{V_k}{A_kP_i}.
\end{equation}

For a spherical vesicle,
\begin{align}
    V_k &= \frac{4}{3}\pi R_k^3, \\
    A_k &= 4\pi R_k^2,
\end{align}
so
\begin{equation}
    \tau_{i,\mathrm{leak}}
    =
    \frac{R_k}{3P_i}.
\end{equation}

The trapped enrichment survives if:
\begin{equation}
    \tau_{i,\mathrm{leak}}
    >
    \tau_{\mathrm{wet}},
\end{equation}
or, more generally,
\begin{equation}
    \tau_{i,\mathrm{leak}}
    >
    \tau_{\mathrm{react,int}}.
\end{equation}

\subsection{Central timescale window}

The wet-dry mechanism works only inside a timescale window:
\begin{equation}
    \tau_{\mathrm{synth},L},\,
    \tau_{\mathrm{asm}}
    <
    \tau_{\mathrm{dry}}
    <
    \tau_{\mathrm{leak}}.
\end{equation}

If the dry phase is too short, amphiphiles do not form or assemble. If the wet
phase is too long relative to leakage, the trapped chemical memory is erased.

\section{Compartment Classes}

Pores, coacervates, and amphiphile vesicles must be treated as different
objects. They are all compartments, but they have different boundary physics.

\subsection{Generic compartment}

Define a generic compartment $k$ as:
\begin{equation}
    \mathcal{C}_k(t)
    =
    \left[
    \Omega_k(t),
    \partial\Omega_k(t),
    r_k(t),
    V_k(t),
    \{X_i^k(t)\},
    B_k
    \right].
\end{equation}

Here $\Omega_k$ is the interior region, $\partial\Omega_k$ is the boundary,
$r_k$ is the position, $V_k$ is the volume, $X_i^k$ is the molecule number of
species $i$, and $B_k$ is the boundary rule.

\subsection{Mineral pore compartment}

A pore is a fixed environmental compartment. It is not a free protocell.

\begin{equation}
    \mathcal{C}_j^{\mathrm{pore}}
    =
    \left[
    \Omega_j^{\mathrm{pore}},
    \partial\Omega_j^{\mathrm{mineral}},
    V_j,
    \{c_i^j(t)\},
    \{S_i^{\mathrm{wall}}\},
    \{K_{jm}^{(i)}\}
    \right].
\end{equation}

The pore equation is:
\begin{equation}
    \frac{dc_i^j}{dt}
    =
    R_i^{\mathrm{aq}}(\mathbf{c}^j;T_j,\mathrm{pH}_j,\mathrm{redox}_j)
    +
    S_i^{\mathrm{wall}}
    +
    \sum_m K_{jm}^{(i)}(c_i^m-c_i^j)
    +
    S_i^{\mathrm{geo}}
    -
    \frac{c_i^j}{h_j}\frac{dh_j}{dt}.
\end{equation}

The boundary rule is:
\begin{equation}
    \partial\Omega_j^{\mathrm{pore}} =
    \mathrm{mineral\ wall}.
\end{equation}

Pores are initial geometry. They do not divide and they are not born from the
chemistry.

\subsection{Coacervate droplet compartment}

A coacervate is a mobile liquid-liquid phase separated droplet. It is a weak
protocell candidate, but it is not a membrane-bounded protocell.

\begin{equation}
    \mathcal{C}_k^{\mathrm{coa}}
    =
    \left[
    r_k,
    R_k,
    V_k,
    \phi_k,
    \{X_i^k\},
    \{K_i^k\},
    \{\tau_{i,\mathrm{part}}^k\}
    \right].
\end{equation}

Here $\phi_k$ is the dense-phase volume fraction, $K_i^k$ is the partition
coefficient, and $\tau_{i,\mathrm{part}}^k$ is the exchange/partitioning time.

The boundary rule is:
\begin{equation}
    \partial\Omega_k^{\mathrm{coa}}
    =
    \mathrm{liquid-liquid\ interface}.
\end{equation}

The partition coefficient is:
\begin{equation}
    K_i^k =
    \frac{c_i^{k,\mathrm{in}}}{c_i^{\mathrm{ext}}}.
\end{equation}

If $K_i^k>1$, species $i$ is enriched inside the coacervate. If $K_i^k<1$,
it is excluded.

Define a local phase-separation score:
\begin{equation}
    \Phi_{\mathrm{coa}}(r,t)
    =
    \sum_{a\in A}
    \sum_{b\in B}
    \alpha_{ab}c_a(r,t)c_b(r,t),
\end{equation}
where $A$ and $B$ are phase-separating species classes, for example oppositely
charged oligomers.

Coacervate formation is allowed when:
\begin{equation}
    \Phi_{\mathrm{coa}}(r,t) >
    \Phi_{\mathrm{coa}}^{\mathrm{crit}}.
\end{equation}

At formation:
\begin{equation}
    X_i^k(t_{\mathrm{form}})
    \sim
    \mathrm{Poisson}
    \left[
    K_i^k c_i^{\mathrm{ext}}(r_k,t_{\mathrm{form}})V_k
    \right].
\end{equation}

After formation:
\begin{equation}
    \frac{dc_i^{k,\mathrm{in}}}{dt}
    =
    R_i^k(\mathbf{c}^k;T,\mathrm{pH},\mathrm{redox})
    +
    \frac{
    K_i^k c_i^{\mathrm{ext}}(r_k,t)
    -
    c_i^{k,\mathrm{in}}
    }
    {\tau_{i,\mathrm{part}}^k}.
\end{equation}

A coacervate retains memory of a dry phase only if:
\begin{equation}
    \tau_{i,\mathrm{part}}^k >
    \tau_{\mathrm{wet}}.
\end{equation}

\subsection{Amphiphile vesicle protocell}

An amphiphile vesicle is the clean object to call a protocell. It is a mobile
membrane-bounded compartment with finite permeability.

\begin{equation}
    \mathcal{P}_k^{\mathrm{ves}}
    =
    \left[
    r_k,
    R_k,
    V_k,
    A_k,
    N_L^k,
    \{X_i^k\},
    \{P_i^k\}
    \right].
\end{equation}

The boundary rule is:
\begin{equation}
    \partial\Omega_k^{\mathrm{ves}}
    =
    \mathrm{amphiphile\ membrane}.
\end{equation}

The membrane exchange equation is:
\begin{equation}
    J_i^k =
    A_kP_i^k
    \left[
    c_i^{\mathrm{ext}}(r_k,t)
    -
    c_i^{k,\mathrm{int}}(t)
    \right],
\end{equation}
with
\begin{equation}
    c_i^{k,\mathrm{int}} = \frac{X_i^k}{V_k}.
\end{equation}

For initial radius $R_0$, required amphiphiles are:
\begin{equation}
    N_{L,\mathrm{req}} =
    \frac{4\pi R_0^2}{a_L}
\end{equation}
for a monolayer or interface coating, and approximately
\begin{equation}
    N_{L,\mathrm{req}} =
    \frac{8\pi R_0^2}{a_L}
\end{equation}
for a bilayer vesicle.

The membrane coverage is:
\begin{equation}
    \theta_k =
    \frac{N_L^k}{N_{L,\mathrm{req}}^k}.
\end{equation}

The vesicle is stable if:
\begin{equation}
    \theta_k \geq \theta_{\min}.
\end{equation}

If $\theta_k<\theta_{\min}$, the vesicle is leaky or dissolves.

\subsection{Coacervate to vesicle transition}

A coacervate can become amphiphile-stabilized by accumulating amphiphiles at
its interface:
\begin{equation}
    \frac{dN_L^k}{dt}
    =
    k_{\mathrm{ads}}A_kc_L^{\mathrm{ext}}
    -
    k_{\mathrm{des}}N_L^k
    +
    R_L^k.
\end{equation}

Here $R_L^k$ is internal or interfacial production of amphiphiles. When
\begin{equation}
    \theta_k =
    \frac{N_L^k}{N_{L,\mathrm{req}}^k}
    \geq 1,
\end{equation}
the transition is:
\begin{equation}
    \mathcal{C}_k^{\mathrm{coa}}
    \rightarrow
    \mathcal{P}_k^{\mathrm{ves}}.
\end{equation}

\subsection{Comparison}

\begin{center}
\begin{tabular}{llll}
\toprule
Object & Boundary & Main rule & Interpretation \\
\midrule
Pore & mineral wall & exchange + wall catalysis & environmental reactor \\
Coacervate & liquid-liquid interface & partitioning $K_i$ & weak protocell \\
Vesicle & amphiphile membrane & permeability $P_i$ & membrane protocell \\
\bottomrule
\end{tabular}
\end{center}

\section{Phase 1: Porous Mineral or Ice Compartments}

\subsection{Physical picture}

Before true membranes exist, compartments can be provided by pores, cracks,
hydrothermal precipitates, pumice-like volcanic rock, clays, or ice eutectic
channels. These structures do not require amphiphiles. They concentrate
reactants and keep reaction products near catalytic surfaces.

This phase is important because it gives a physical route to:
\begin{enumerate}
    \item local concentration;
    \item surface catalysis;
    \item partial isolation;
    \item reaction-diffusion gradients;
    \item repeated feeding and flushing.
\end{enumerate}

\subsection{Pore concentration}

Let a pore have radius $r_p$, length $L_p$, volume $V_p$, and surface area
$A_p$:
\begin{align}
    V_p &= \pi r_p^2 L_p, \\
    A_p &= 2\pi r_p L_p.
\end{align}

The surface-to-volume ratio is
\begin{equation}
    \frac{A_p}{V_p} = \frac{2}{r_p}.
\end{equation}

Small pores strongly enhance surface chemistry because $A_p/V_p$ increases as
$r_p$ decreases.

If a dissolved species $i$ is concentrated by evaporation, freezing, or
adsorption, write:
\begin{equation}
    C_{i,\mathrm{pore}} =
    \Gamma_{\mathrm{conc}} C_{i,\mathrm{bulk}},
\end{equation}
where
\begin{equation}
    \Gamma_{\mathrm{conc}}
    =
    \Gamma_{\mathrm{wetdry}}
    \Gamma_{\mathrm{freeze}}
    \Gamma_{\mathrm{ads}}
    \Gamma_{\mathrm{flow}}.
\end{equation}

\subsection{Diffusion and residence time}

For a species with molecular diffusion coefficient $D_i$, the diffusion time
across a pore of length $L_p$ is
\begin{equation}
    t_{\mathrm{diff},i}
    =
    \frac{L_p^2}{D_i}.
\end{equation}

If there is advective flow velocity $u$, the advective residence time is
\begin{equation}
    t_{\mathrm{adv}} =
    \frac{L_p}{u}.
\end{equation}

The Damkohler number compares reaction and transport:
\begin{equation}
    \mathrm{Da}_i =
    \frac{t_{\mathrm{transport}}}{t_{\mathrm{reaction},i}}.
\end{equation}

For a first-order reaction with rate constant $k_i$:
\begin{equation}
    t_{\mathrm{reaction},i} = \frac{1}{k_i}.
\end{equation}

Chemistry is efficient inside pores when
\begin{equation}
    \mathrm{Da}_i \gtrsim 1.
\end{equation}

\subsection{Surface-catalyzed rate}

For a mineral-catalyzed reaction
\begin{equation}
    A + B \rightarrow C,
\end{equation}
use:
\begin{equation}
    R_C =
    k_T f_{\mathrm{pH}} f_I
    \left(1+\alpha_m S_m\right)
    C_A C_B.
\end{equation}

The temperature factor is:
\begin{equation}
    k_T = A_0 \exp\left(-\frac{E_a}{RT}\right).
\end{equation}

The pH window can be represented by:
\begin{equation}
    f_{\mathrm{pH}}
    =
    \exp\left[-\left(\frac{\mathrm{pH}-\mathrm{pH}_{\mathrm{opt}}}
    {\sigma_{\mathrm{pH}}}\right)^2\right].
\end{equation}

\subsection{Pore compartment viability criterion}

A porous compartment is chemically useful if:
\begin{align}
    \Gamma_{\mathrm{conc}} &> 1, \\
    \mathrm{Da}_{\mathrm{key}} &\gtrsim 1, \\
    S_m(X) &> S_{\mathrm{min}}, \\
    t_{\mathrm{res}} &> t_{\mathrm{polymerization}}, \\
    T_{\mathrm{env}} &\in [T_{\mathrm{min}},T_{\mathrm{max}}].
\end{align}

This gives a computable Boolean:
\begin{equation}
    \mathcal{P}_{\mathrm{pore}} =
    \Theta(\Gamma_{\mathrm{conc}}-1)
    \Theta(\mathrm{Da}_{\mathrm{key}}-1)
    \Theta(S_m-S_{\mathrm{min}})
    \Theta(t_{\mathrm{res}}-t_{\mathrm{poly}}).
\end{equation}

\section{Phase 2: Solid Compartments and Surface-Bound Networks}

\subsection{Physical picture}

Solid compartments are not closed cells. They are reaction domains where
molecules are adsorbed on mineral or ice surfaces. Examples include clay
interlayers, FeS/NiS hydrothermal precipitates, phosphate-rich surfaces, silica
surfaces, volcanic rock pores, and ice grain boundaries.

This phase can support reaction networks before membranes exist.

\subsection{Adsorption equilibrium}

For adsorption of species $i$ onto a surface $S$:
\begin{equation}
    i + S \rightleftharpoons iS.
\end{equation}

A Langmuir model gives:
\begin{equation}
    \theta_i =
    \frac{K_i C_i}{1+\sum_j K_j C_j},
\end{equation}
where $\theta_i$ is fractional surface coverage.

The adsorbed concentration is:
\begin{equation}
    C_{i,S} =
    \theta_i S_{\mathrm{site}},
\end{equation}
where $S_{\mathrm{site}}$ is the concentration of reactive surface sites.

\subsection{Surface polymerization}

For monomer addition on a surface:
\begin{equation}
    P_nS + M \rightarrow P_{n+1}S,
\end{equation}
use:
\begin{equation}
    \frac{dC_{P_{n+1},S}}{dt}
    =
    k_{\mathrm{surf}}(T,\mathrm{pH},I)
    C_{P_n,S} C_M
    -
    k_{\mathrm{hyd}}(T,\mathrm{pH})
    C_{P_{n+1},S}.
\end{equation}

For peptide-like growth:
\begin{equation}
    \mathrm{pep}_n + \mathrm{AA}^{*}
    \rightarrow
    \mathrm{pep}_{n+1} + \mathrm{leaving\ group},
\end{equation}
where $\mathrm{AA}^{*}$ is an activated amino acid such as an aminoacyl
thioester or aminoacyl phosphate.

\subsection{Surface retention versus loss}

Surface compartments matter only if useful products are retained long enough.
For adsorbed polymer $P_n$:
\begin{equation}
    \frac{dC_{P_n,S}}{dt}
    =
    R_{\mathrm{form}}
    -
    R_{\mathrm{hyd}}
    -
    k_{\mathrm{des}} C_{P_n,S}
    -
    k_{\mathrm{flush}} C_{P_n,S}.
\end{equation}

The retention condition is:
\begin{equation}
    k_{\mathrm{form}}
    >
    k_{\mathrm{hyd}} + k_{\mathrm{des}} + k_{\mathrm{flush}}.
\end{equation}

\section{Phase 3: Coacervate Compartments}

\subsection{Physical picture}

Coacervates are liquid-liquid phase separated droplets. They can form from
oppositely charged polymers, peptides, nucleotides, polyions, or mixed organic
polymers. They are not lipid membranes, but they can concentrate solutes and
create chemically distinct microenvironments.

For LIFEORIG, coacervates are useful as an intermediate between surface-bound
chemistry and amphiphilic vesicles.

\subsection{Phase separation criterion}

Let $\phi$ be the polymer volume fraction and $\chi$ an effective interaction
parameter. A simple Flory-Huggins free energy density is:
\begin{equation}
    \frac{\Delta G_{\mathrm{mix}}}{k_BT}
    =
    \frac{\phi}{N}\ln\phi
    +
    (1-\phi)\ln(1-\phi)
    +
    \chi \phi(1-\phi),
\end{equation}
where $N$ is the polymerization degree.

Phase separation is possible when:
\begin{equation}
    \frac{\partial^2 \Delta G_{\mathrm{mix}}}{\partial \phi^2} < 0.
\end{equation}

The spinodal condition is:
\begin{equation}
    \frac{1}{N\phi} + \frac{1}{1-\phi} - 2\chi < 0.
\end{equation}

Therefore a rough threshold is:
\begin{equation}
    \chi >
    \frac{1}{2}
    \left[
    \frac{1}{N\phi} + \frac{1}{1-\phi}
    \right].
\end{equation}

\subsection{Salt and pH dependence}

Coacervation depends strongly on ionic strength and charge state. Represent
the effective interaction parameter as:
\begin{equation}
    \chi_{\mathrm{eff}}
    =
    \chi_0
    +
    \chi_{\mathrm{charge}}(\mathrm{pH})
    -
    \chi_I I.
\end{equation}

Polymer charge can be approximated from acid-base groups:
\begin{equation}
    q(\mathrm{pH})
    =
    \sum_a
    \frac{-1}{1+10^{pK_{a,a}-\mathrm{pH}}}
    +
    \sum_b
    \frac{+1}{1+10^{\mathrm{pH}-pK_{a,b}}}.
\end{equation}

Coacervate stability is favored when oppositely charged species coexist:
\begin{equation}
    q_1(\mathrm{pH}) q_2(\mathrm{pH}) < 0.
\end{equation}

\subsection{Partitioning into droplets}

For species $i$, define a partition coefficient:
\begin{equation}
    K_i^{\mathrm{coac}}
    =
    \frac{C_{i,\mathrm{droplet}}}{C_{i,\mathrm{bulk}}}.
\end{equation}

Reaction rates inside coacervates use:
\begin{equation}
    C_{i,\mathrm{local}} =
    K_i^{\mathrm{coac}} C_{i,\mathrm{bulk}}.
\end{equation}

Thus, for a reaction $A+B\rightarrow C$:
\begin{equation}
    R_{C,\mathrm{coac}}
    =
    k(T,\mathrm{pH},I)
    K_A^{\mathrm{coac}} K_B^{\mathrm{coac}}
    C_{A,\mathrm{bulk}} C_{B,\mathrm{bulk}}.
\end{equation}

\subsection{Droplet growth}

Let $R_d$ be droplet radius. If growth is controlled by uptake of material
from solution:
\begin{equation}
    \frac{dR_d}{dt}
    =
    \frac{D}{R_d}
    \left(
    \frac{C_{\infty}-C_{\mathrm{eq}}}{\rho_d}
    \right).
\end{equation}

Coarsening can be represented by a Lifshitz-Slyozov-like law:
\begin{equation}
    R_d^3(t) - R_d^3(0) = K_{\mathrm{LSW}} t.
\end{equation}

\subsection{Coacervate viability criterion}

Coacervates are viable chemical compartments if:
\begin{align}
    \chi_{\mathrm{eff}} &> \chi_{\mathrm{crit}}, \\
    K_i^{\mathrm{coac}} &> 1
    \quad \mathrm{for\ key\ feedstocks/catalysts}, \\
    t_{\mathrm{life,droplet}} &> t_{\mathrm{reaction}}, \\
    k_{\mathrm{leak}} &< k_{\mathrm{production}}
    \quad \mathrm{for\ retained\ products}.
\end{align}

\section{Phase 4: Amphiphilic Vesicles and True Protocells}

\subsection{Physical picture}

An amphiphilic protocell is a compartment bounded by a membrane-like assembly:
fatty acid vesicles, mixed amphiphile vesicles, lipid-like membranes, or
related structures. This is closer to a true protocell because it provides
encapsulation, selective permeability, growth, division, and inheritance of
internal chemistry.

The key transition is:
\begin{equation}
    \mathrm{compartment}
    +
    \mathrm{chemical\ network}
    +
    \mathrm{growth/division}
    \rightarrow
    \mathrm{protocell}.
\end{equation}

\subsection{Amphiphile production threshold}

Let $C_A$ be total amphiphile concentration. Vesicles require:
\begin{equation}
    C_A > \mathrm{CMC},
\end{equation}
where CMC is the critical micelle/aggregate concentration.

The CMC depends on temperature, pH, ionic strength, and amphiphile chain
length:
\begin{equation}
    \mathrm{CMC}
    =
    \mathrm{CMC}_0
    \exp\left[
    \frac{\Delta G_{\mathrm{agg}}(T,\mathrm{pH},I,N_c)}
    {RT}
    \right].
\end{equation}

Here $N_c$ is the hydrocarbon chain length or an equivalent hydrophobicity
index.

\subsection{Vesicle geometry}

For a spherical vesicle of radius $R_v$:
\begin{align}
    V_v &= \frac{4}{3}\pi R_v^3, \\
    A_v &= 4\pi R_v^2.
\end{align}

If each amphiphile occupies membrane area $a_0$, the number of amphiphiles in
one leaflet is:
\begin{equation}
    N_{\mathrm{leaflet}} =
    \frac{4\pi R_v^2}{a_0}.
\end{equation}

A bilayer vesicle has approximately:
\begin{equation}
    N_{\mathrm{amp}} \approx
    \frac{8\pi R_v^2}{a_0}.
\end{equation}

\subsection{Encapsulation statistics}

If a solute has concentration $C_i$, the expected number of molecules inside a
vesicle is:
\begin{equation}
    \lambda_i = C_i N_A V_v.
\end{equation}

The probability that the vesicle contains at least one molecule of $i$ is:
\begin{equation}
    P(N_i\geq 1) = 1-\exp(-\lambda_i).
\end{equation}

For an autocatalytic set requiring species $i=1,\ldots,m$, assuming independent
encapsulation:
\begin{equation}
    P_{\mathrm{set}}
    =
    \prod_{i=1}^{m}
    \left[1-\exp(-C_i N_A V_v)\right].
\end{equation}

This connects directly to the idea that autocatalysis requires a sufficient
number of molecules. If $C_i N_A V_v \ll 1$, most vesicles do not contain the
required species.

\subsection{Membrane permeability}

For solute $i$, membrane flux is:
\begin{equation}
    J_i =
    P_i \left(C_{i,\mathrm{out}}-C_{i,\mathrm{in}}\right),
\end{equation}
where $P_i$ is permeability. The concentration change inside a vesicle is:
\begin{equation}
    \frac{dC_{i,\mathrm{in}}}{dt}
    =
    \frac{A_v}{V_v}
    P_i
    \left(C_{i,\mathrm{out}}-C_{i,\mathrm{in}}\right)
    +
    R_{i,\mathrm{internal}}.
\end{equation}

Because
\begin{equation}
    \frac{A_v}{V_v} = \frac{3}{R_v},
\end{equation}
small vesicles exchange molecules faster than large vesicles.

\subsection{Osmotic pressure}

The osmotic pressure difference is:
\begin{equation}
    \Delta \Pi =
    RT
    \sum_i
    \left(
    C_{i,\mathrm{in}} - C_{i,\mathrm{out}}
    \right).
\end{equation}

Membrane tension due to pressure difference can be estimated by Laplace's law:
\begin{equation}
    \Delta P =
    \frac{2\gamma}{R_v},
\end{equation}
where $\gamma$ is membrane tension. Vesicle rupture or division can occur when
$\gamma$ exceeds a critical tension $\gamma_c$.

\subsection{Growth}

If amphiphiles are supplied from the environment, membrane area grows as:
\begin{equation}
    \frac{dA_v}{dt}
    =
    a_0 J_A A_v,
\end{equation}
where $J_A$ is the net amphiphile insertion flux per unit membrane area.

For spherical vesicles:
\begin{equation}
    \frac{dR_v}{dt}
    =
    \frac{a_0 J_A R_v}{2}.
\end{equation}

Internal chemistry grows if:
\begin{equation}
    \frac{dN_{\mathrm{cat}}}{dt}
    =
    R_{\mathrm{cat}} V_v
    -
    k_{\mathrm{loss}}N_{\mathrm{cat}}
    >
    0.
\end{equation}

\subsection{Division}

A simple division criterion can be based on surface-to-volume imbalance. If
membrane area grows faster than volume:
\begin{equation}
    \frac{dA_v/A_v}{dt}
    >
    \frac{2}{3}
    \frac{dV_v/V_v}{dt},
\end{equation}
the vesicle becomes geometrically unstable.

A phenomenological division rate can be written:
\begin{equation}
    k_{\mathrm{div}}
    =
    k_0
    \Theta(A_v-A_{\mathrm{crit}})
    \Theta(\gamma-\gamma_{\mathrm{crit}}).
\end{equation}

\section{Autocatalysis and Catalysis Probability}

\subsection{Molecule-number criterion}

For a catalytic polymer species $P$, the expected copy number in a compartment
of volume $V_c$ is:
\begin{equation}
    \lambda_P = C_P N_A V_c.
\end{equation}

The probability that the compartment contains at least $N_{\mathrm{min}}$
copies is:
\begin{equation}
    P(N_P \geq N_{\mathrm{min}})
    =
    1 -
    \sum_{n=0}^{N_{\mathrm{min}}-1}
    \frac{\lambda_P^n e^{-\lambda_P}}{n!}.
\end{equation}

\subsection{Catalysis probability}

Let a polymer of length $n$ have probability
\begin{equation}
    p_{\mathrm{cat}}(n,c)
\end{equation}
of catalyzing reaction class $c$. A first parameterization is:
\begin{equation}
    p_{\mathrm{cat}}(n,c)
    =
    p_{0,c}
    \left[
    1-\exp\left(-\frac{n-n_{\min,c}}{n_{0,c}}\right)
    \right]
    \Theta(n-n_{\min,c}).
\end{equation}

The effective catalyzed rate for reaction $r$ is:
\begin{equation}
    R_r =
    R_{r,0}
    \left[
    1 +
    \sum_P
    \alpha_{rP}
    p_{\mathrm{cat}}(n_P,c_r)
    N_P
    \right].
\end{equation}

\subsection{Autocatalytic closure criterion}

A reaction set is autocatalytic if there exists a subset $\mathcal{S}$ such
that:
\begin{enumerate}
    \item every reaction in $\mathcal{S}$ is catalyzed by at least one molecule
    produced by $\mathcal{S}$;
    \item every reactant can be made from the food set plus products of
    $\mathcal{S}$.
\end{enumerate}

Computationally, define a catalysis matrix:
\begin{equation}
    K_{ij} =
    \begin{cases}
    1, & \mathrm{species\ } i \mathrm{\ catalyzes\ reaction\ } j, \\
    0, & \mathrm{otherwise}.
    \end{cases}
\end{equation}

Define stoichiometry matrix $\nu_{ij}$:
\begin{equation}
    \frac{dC_i}{dt}
    =
    \sum_j \nu_{ij} R_j(C,T,\mathrm{pH},I).
\end{equation}

Autocatalytic growth is detected when a positive eigenvalue appears in the
linearized production operator around the feedstock state:
\begin{equation}
    \max \mathrm{Re}\left[\lambda
    \left(
    \frac{\partial \dot{C}}{\partial C}
    \right)\right] > 0.
\end{equation}

\section{Transitions Between Phases}

\subsection{Pores to coacervates}

The transition from pore chemistry to coacervates requires sufficient polymer
production and phase separation:
\begin{align}
    C_{\mathrm{poly}} &> C_{\mathrm{LLPS}}, \\
    \chi_{\mathrm{eff}} &> \chi_{\mathrm{crit}}, \\
    K_i^{\mathrm{coac}} &> 1
    \quad \mathrm{for\ catalytic\ species}.
\end{align}

\subsection{Pores or coacervates to amphiphilic vesicles}

The transition to vesicles requires amphiphile production:
\begin{align}
    C_A &> \mathrm{CMC}, \\
    T_{\mathrm{env}} &\in [T_{\mathrm{mem,min}},T_{\mathrm{mem,max}}], \\
    \mathrm{pH} &\in [\mathrm{pH}_{\mathrm{mem,min}},
                     \mathrm{pH}_{\mathrm{mem,max}}], \\
    I &< I_{\mathrm{destabilize}}.
\end{align}

\subsection{Vesicle to protocell}

A vesicle becomes protocell-like when it couples internal chemistry to
compartment growth:
\begin{align}
    P_{\mathrm{set}} &> P_{\mathrm{min}}, \\
    \frac{dN_{\mathrm{cat}}}{dt} &> 0, \\
    \frac{dA_v}{dt} &> 0, \\
    k_{\mathrm{div}} &> 0, \\
    \mathrm{heritability} &> h_{\mathrm{min}}.
\end{align}

Here heritability can be defined operationally as the correlation between
parent and daughter catalytic composition:
\begin{equation}
    h =
    \mathrm{corr}
    \left(
    \mathbf{N}_{\mathrm{cat,parent}},
    \mathbf{N}_{\mathrm{cat,daughter}}
    \right).
\end{equation}

\section{Minimal Dynamical Model}

For each compartment type $c$, evolve:
\begin{align}
    \frac{dC_i^{(c)}}{dt}
    &=
    \sum_j \nu_{ij} R_j^{(c)}
    +
    \frac{A_c}{V_c}
    P_i^{(c)}
    \left(C_i^{\mathrm{env}}-C_i^{(c)}\right)
    -
    k_{\mathrm{loss},i}^{(c)} C_i^{(c)}, \\
    \frac{dV_c}{dt}
    &=
    G_c(C,T,\mathrm{pH},I,a_w), \\
    \frac{dN_c}{dt}
    &=
    k_{\mathrm{birth}}^{(c)}N_c
    -
    k_{\mathrm{death}}^{(c)}N_c.
\end{align}

For pores, $V_c$ is usually fixed. For coacervates, $V_c$ changes by phase
separation and coarsening. For vesicles, $V_c$ and membrane area evolve through
amphiphile uptake, osmotic pressure, and division.

\section{Implementation Pathway}

\subsection{Step 1: define a ProtocellEnvironmentState}

The state should contain:
\begin{align}
    T,\ P,\ \mathrm{pH},\ I,\ a_w,\ F_{\mathrm{UV}},
    \Phi_{\mathrm{wetdry}},\ \Phi_{\mathrm{freeze}},
    S_m(X),\ C_i.
\end{align}

\subsection{Step 2: evaluate compartment availability}

Compute:
\begin{align}
    \mathcal{P}_{\mathrm{pore}},\quad
    \mathcal{P}_{\mathrm{surface}},\quad
    \mathcal{P}_{\mathrm{coac}},\quad
    \mathcal{P}_{\mathrm{vesicle}}.
\end{align}

These are not final probabilities at first. They can start as Boolean
viability flags.

\subsection{Step 3: run chemistry inside each viable compartment}

For each viable compartment:
\begin{equation}
    \frac{d\mathbf{C}^{(c)}}{dt}
    =
    \mathbf{S}\mathbf{R}^{(c)}
    +
    \mathbf{F}_{\mathrm{exchange}}^{(c)}
    -
    \mathbf{L}^{(c)}.
\end{equation}

\subsection{Step 4: test autocatalytic threshold}

Track:
\begin{align}
    N_{\mathrm{polymer}},\quad
    P_{\mathrm{set}},\quad
    p_{\mathrm{cat}},\quad
    \max \mathrm{Re}(\lambda).
\end{align}

The first computational success condition is not ``life''. It is:
\begin{equation}
    \mathrm{persistent\ compartmentalized\ autocatalytic\ chemistry}.
\end{equation}

\section{Reference Starting Points}

The following literature is useful for model construction:
\begin{itemize}
    \item Deamer and Dworkin: chemistry and physics of amphiphilic
    compartments and prebiotic membranes.
    \item Szostak and collaborators: fatty-acid vesicles, protocell growth,
    membrane permeability, and encapsulation.
    \item Hanczyc, Fujikawa, and Szostak: experimental growth and division of
    fatty-acid vesicles.
    \item Monnard and Deamer: membrane self-assembly under prebiotic
    conditions.
    \item Budin and Szostak: physical effects of membrane composition and
    competition between protocells.
    \item Koga, Williams, Perriman, and Mann: protocell models based on
    coacervates and phase-separated droplets.
    \item Dora Tang and collaborators: coacervates, protocell droplets, and
    biochemical compartmentalization.
    \item Martin, Russell, and collaborators: alkaline hydrothermal vents,
    mineral compartments, and geochemical energy gradients.
    \item Kauffman: autocatalytic sets and collective autocatalysis.
    \item Hordijk and Steel: RAF theory for autocatalytic reaction networks.
\end{itemize}

\section{Working Interpretation}

The sequence to test in LIFEORIG should be:
\begin{equation}
    \mathrm{pores/surfaces}
    \rightarrow
    \mathrm{polymer\ enrichment}
    \rightarrow
    \mathrm{coacervate\ or\ vesicle\ compartment}
    \rightarrow
    \mathrm{autocatalytic\ chemistry}
    \rightarrow
    \mathrm{growth/division}.
\end{equation}

For surface volcanic environments, wet/dry pore and vesicle pathways are
probably most relevant. For hydrothermal vents, mineral pore and surface-bound
networks are the natural first compartment. For icy ocean worlds, eutectic ice
channels and hydrothermal pore compartments may be more plausible than surface
pond vesicles, unless there is local exposure, brine cycling, or impact/thermal
processing.

\end{document}
